% ============================================================================
% CHAPTER 1: CLASSICAL NLP FOUNDATIONS
% ============================================================================

\chapter{Classical NLP Foundations}
\label{chap:classical_foundations}

\begin{tldr}
Classical NLP foundations---text preprocessing, representation, and retrieval---remain
essential for production systems and interviews alike. This chapter covers the full text
processing pipeline, classical text representations (Bag of Words, TF-IDF, BM25), syntactic
analysis fundamentals, and the core NLP task taxonomy. Understanding these foundations is
critical because many production systems still rely on classical methods for their efficiency,
interpretability, and robustness, and interviewers test them to assess whether a candidate has
genuine depth or is merely an ``LLM tourist'' who cannot reason below the API layer.
\end{tldr}

% ============================================================================
\section{The NLP Landscape}
\label{sec:nlp_landscape}
% ============================================================================

Natural Language Processing is the subfield of artificial intelligence concerned with enabling
computers to understand, generate, and reason about human language. Although modern large
language models have captured public attention, the discipline spans decades of research
across linguistics, statistics, and computer science. A well-prepared candidate must
understand not only \emph{what} the current state of the art is, but \emph{how} and
\emph{why} the field arrived here.

\subsection{Why NLP Is Hard}
\label{sec:why_nlp_hard}

Language presents unique challenges that distinguish NLP from other machine learning
domains. Interviewers frequently probe whether a candidate appreciates these difficulties
or naively assumes that ``more data and bigger models'' solve everything.

\begin{enumerate}[label=\textbf{\arabic*.}]
    \item \textbf{Ambiguity.} Natural language is ambiguous at every level:
    \begin{itemize}
        \item \emph{Lexical ambiguity:} ``bank'' can mean a financial institution or a river
              bank.
        \item \emph{Syntactic ambiguity:} ``I saw the man with the telescope''---who has the
              telescope?
        \item \emph{Semantic ambiguity:} ``Every student read a book''---the same book or
              different books?
        \item \emph{Pragmatic ambiguity:} ``Can you pass the salt?'' is a request, not a
              question about ability.
    \end{itemize}

    \item \textbf{Context Dependence.} The meaning of a word, phrase, or sentence depends
    heavily on surrounding context. Pronouns require coreference resolution (``John told
    Mary he would leave''---who is ``he''?). Discourse-level phenomena like topic shifts,
    implicature, and presupposition make isolated sentence processing insufficient.

    \item \textbf{Compositionality.} The meaning of a sentence is built from the meanings of
    its parts and the rules used to combine them. ``The dog bit the man'' versus ``The man
    bit the dog''---same words, opposite meanings. Models must capture this compositional
    structure.

    \item \textbf{Data Sparsity.} The space of possible sentences is effectively infinite.
    Zipf's law tells us that word frequency follows a power law: a small number of words
    account for most occurrences, while the vast majority of words are rare. This creates
    severe challenges for count-based methods.

    \item \textbf{Multilinguality.} There are approximately 7,000 human languages. They
    differ in script (Latin, Cyrillic, Devanagari, CJK), morphology (English is
    relatively analytic; Turkish, Finnish, and Korean are highly agglutinative), word
    order (SVO, SOV, VSO), and writing conventions (right-to-left, no whitespace
    between words). Any NLP system deployed globally must contend with this diversity.

    \item \textbf{World Knowledge.} Understanding language often requires knowledge that is
    not present in the text. ``The trophy doesn't fit in the suitcase because it is too
    big''---resolving ``it'' requires physical reasoning about objects and sizes.
\end{enumerate}

\begin{keyinsight}[Why Interviewers Ask About These Challenges]
Asking ``why is NLP hard?'' is a calibration question. A shallow answer lists one or two
points. A strong answer demonstrates understanding of multiple levels of ambiguity, gives
concrete examples, and connects these challenges to why specific architectural choices
(e.g., contextual embeddings, attention mechanisms) were necessary.
\end{keyinsight}

\subsection{The Four Paradigm Shifts}
\label{sec:paradigm_shifts}

The history of NLP can be organized into four major paradigm shifts. Understanding this
progression is essential because interviewers test whether you know \emph{why} each
transition occurred, not just \emph{that} it occurred.

\begin{figure}[H]
\centering
\begin{tikzpicture}[
    era/.style={
        draw=deepblue, fill=lightblue, rounded corners=5pt,
        minimum width=3.8cm, minimum height=1.8cm,
        align=center, font=\small\sffamily
    },
    arrow/.style={-{Stealth[length=3mm]}, thick, deepblue},
    label/.style={font=\scriptsize\sffamily\color{gray!70!black}, align=center}
]
    % Era boxes
    \node[era] (rule) at (0,0)
        {\textbf{Rule-Based}\\[2pt]Hand-crafted rules\\Gazetteers, regex\\CFG parsers};
    \node[era] (stat) at (4.5,0)
        {\textbf{Statistical}\\[2pt]N-gram LMs, HMMs\\CRFs, MaxEnt\\TF-IDF, BM25};
    \node[era] (neural) at (9,0)
        {\textbf{Neural}\\[2pt]Word2Vec, RNNs\\LSTMs, Seq2Seq\\Attention};
    \node[era] (transformer) at (13.5,0)
        {\textbf{Transformer}\\[2pt]BERT, GPT, T5\\LLMs, RLHF\\RAG, Agents};

    % Arrows
    \draw[arrow] (rule) -- (stat);
    \draw[arrow] (stat) -- (neural);
    \draw[arrow] (neural) -- (transformer);

    % Time labels
    \node[label] at (0,-1.4) {1950s--1990s};
    \node[label] at (4.5,-1.4) {1990s--2013};
    \node[label] at (9,-1.4) {2013--2017};
    \node[label] at (13.5,-1.4) {2017--Present};

    % Driving forces (above)
    \node[label] at (2.25,1.3) {Data availability\\Machine learning};
    \node[label] at (6.75,1.3) {GPU compute\\Distributed representations};
    \node[label] at (11.25,1.3) {Self-attention\\Scale + pretraining};
\end{tikzpicture}
\caption{The four paradigm shifts in NLP. Each transition was driven by both algorithmic
insight and increased availability of data and compute.}
\label{fig:nlp_paradigm_shifts}
\end{figure}

\begin{enumerate}[label=\textbf{Era \arabic*.}]
    \item \textbf{Rule-Based (1950s--1990s).} Linguists hand-coded grammar rules,
    gazetteers, and pattern-matching systems. Examples include ELIZA (pattern matching for
    dialogue), SHRDLU (natural language understanding in a blocks world), and early
    machine translation systems. These systems were brittle, expensive to build, and
    failed to generalize, but they established the problems NLP would spend decades
    solving.

    \item \textbf{Statistical (1990s--2013).} The field shifted to learning patterns from
    data. N-gram language models, Hidden Markov Models (HMMs) for sequence labeling,
    Maximum Entropy classifiers, Conditional Random Fields (CRFs), and TF-IDF/BM25
    for information retrieval defined this era. Fred Jelinek's famous quip---``Every
    time I fire a linguist, the performance of the speech recognizer goes up''---captures
    the spirit of this shift. \emph{This era's methods are the focus of this chapter.}

    \item \textbf{Neural (2013--2017).} Word2Vec (Mikolov et al., 2013) showed that
    distributed representations capture semantic relationships. Recurrent neural networks
    (RNNs/LSTMs) enabled sequence modeling, and the attention mechanism (Bahdanau et al.,
    2015) solved the information bottleneck in encoder-decoder models. This era introduced
    end-to-end learning and eliminated much manual feature engineering.

    \item \textbf{Transformer and LLM (2017--Present).} ``Attention Is All You Need''
    (Vaswani et al., 2017) introduced the Transformer, replacing recurrence with
    self-attention. BERT (2018) demonstrated the power of bidirectional pretraining. GPT-3
    (2020) revealed emergent in-context learning at scale. RLHF and instruction tuning
    (2022+) made LLMs practically useful. This era is defined by scale, pretraining, and
    adaptation.
\end{enumerate}

\begin{warningbox}
A common interview mistake is dismissing classical methods as ``outdated.'' BM25 is still
the backbone of production search systems at Google, Amazon, and Elasticsearch. TF-IDF
features are used in spam filters handling billions of messages. CRFs remain competitive for
structured prediction in low-latency settings. Interviewers specifically test classical NLP
to filter out candidates who only have surface-level LLM knowledge.
\end{warningbox}

\subsection{Core NLP Task Taxonomy}
\label{sec:task_taxonomy}

NLP tasks can be organized into a taxonomy based on their input-output structure. This
taxonomy is essential because it determines the appropriate model architecture, loss
function, and evaluation metric.

\begin{figure}[H]
\centering
\begin{tikzpicture}[
    root/.style={draw=deepblue, fill=deepblue, text=white, rounded corners=4pt,
                 font=\small\sffamily\bfseries, minimum width=2.4cm, minimum height=0.8cm},
    cat/.style={draw=deepblue!70, fill=lightblue, rounded corners=3pt,
                font=\small\sffamily, minimum width=2.8cm, minimum height=0.7cm, align=center},
    task/.style={draw=gray!50, fill=white, rounded corners=2pt,
                 font=\scriptsize\sffamily, minimum width=2.3cm, minimum height=0.55cm, align=center},
    edge from parent/.style={draw=deepblue!60, thick, -{Stealth[length=2mm]}},
    level 1/.style={sibling distance=3.8cm, level distance=1.4cm},
    level 2/.style={sibling distance=2.5cm, level distance=1.3cm}
]
\node[root] {NLP Tasks}
    child { node[cat] {Text\\Classification}
        child { node[task] {Sentiment} }
        child { node[task] {Topic} }
        child { node[task] {Intent} }
    }
    child { node[cat] {Sequence\\Labeling}
        child { node[task] {NER} }
        child { node[task] {POS} }
        child { node[task] {Chunking} }
    }
    child { node[cat] {Seq-to-Seq}
        child { node[task] {Translation} }
        child { node[task] {Summarization} }
        child { node[task] {QA} }
    }
    child { node[cat] {Structured\\Prediction}
        child { node[task] {Parsing} }
        child { node[task] {Coreference} }
        child { node[task] {Relation Ext.} }
    };
\end{tikzpicture}
\caption{Core NLP task taxonomy organized by input-output structure. Each category
implies different architectural choices and evaluation metrics.}
\label{fig:task_taxonomy}
\end{figure}

\begin{itemize}
    \item \textbf{Text Classification:} Map an entire text to a label (or set of labels).
    Examples: sentiment analysis, spam detection, intent classification, topic
    categorization. Output is a probability distribution over a discrete label set.

    \item \textbf{Sequence Labeling:} Assign a label to each token in a sequence. Examples:
    Named Entity Recognition (NER), Part-of-Speech (POS) tagging, chunking. The
    critical property is that output labels may have dependencies (e.g., \texttt{I-PER}
    cannot follow \texttt{B-LOC}).

    \item \textbf{Sequence-to-Sequence (Seq2Seq):} Map an input sequence to an output
    sequence of potentially different length. Examples: machine translation,
    abstractive summarization, question answering. Requires generation capability.

    \item \textbf{Structured Prediction:} Produce structured outputs such as parse trees,
    knowledge graphs, or coreference clusters. Often involves combinatorial search and
    specialized decoding algorithms (e.g., CKY for parsing, Viterbi for sequence
    labeling).
\end{itemize}

% ============================================================================
\section{The Text Preprocessing Pipeline}
\label{sec:preprocessing}
% ============================================================================

Text preprocessing transforms raw text into a form suitable for downstream NLP models.
While often treated as a trivial engineering detail, preprocessing decisions profoundly
affect model performance. Interviewers at FAANG companies test preprocessing knowledge
because it reveals whether a candidate has actually built NLP systems or only trained models
on clean benchmarks.

\begin{figure}[H]
\centering
\begin{tikzpicture}[
    step/.style={draw=deepblue, fill=lightblue, rounded corners=4pt,
                 font=\small\sffamily, minimum width=2.6cm, minimum height=0.9cm, align=center},
    arrow/.style={-{Stealth[length=2.5mm]}, thick, deepblue!70},
    note/.style={font=\scriptsize\sffamily\itshape, text=gray!70!black, align=center}
]
    \node[step] (raw) at (0,0) {Raw Text};
    \node[step] (enc) at (3.2,0) {Encoding\\Normalization};
    \node[step] (uni) at (6.4,0) {Unicode\\Normalization};
    \node[step] (tok) at (9.6,0) {Tokenization};
    \node[step] (case) at (0,-2.2) {Case\\Normalization};
    \node[step] (stop) at (3.2,-2.2) {Stopword\\Removal};
    \node[step] (stem) at (6.4,-2.2) {Stemming /\\Lemmatization};
    \node[step] (feat) at (9.6,-2.2) {Feature\\Extraction};

    \draw[arrow] (raw) -- (enc);
    \draw[arrow] (enc) -- (uni);
    \draw[arrow] (uni) -- (tok);
    \draw[arrow] (tok) -- ++(0,-0.55) -| (case);
    \draw[arrow] (case) -- (stop);
    \draw[arrow] (stop) -- (stem);
    \draw[arrow] (stem) -- (feat);

    % Notes
    \node[note] at (3.2,0.75) {UTF-8};
    \node[note] at (6.4,0.75) {NFKC};
    \node[note] at (0,-1.45) {Optional};
    \node[note] at (3.2,-1.45) {Task-\\dependent};
    \node[note] at (6.4,-1.45) {Porter / WordNet};
    \node[note] at (9.6,-1.45) {BoW / TF-IDF /\\N-grams};
\end{tikzpicture}
\caption{The classical text preprocessing pipeline. Each step is optional and
task-dependent; blindly applying all steps is a common beginner mistake.}
\label{fig:preprocessing_pipeline}
\end{figure}

\subsection{Tokenization}
\label{sec:tokenization}

Tokenization is the process of splitting text into discrete units (tokens). Despite
appearing trivial, it is one of the most consequential decisions in any NLP pipeline.

\subsubsection{Whitespace and Rule-Based Tokenization}

The simplest approach splits text on whitespace characters. This works reasonably well for
English but fails in numerous cases:

\begin{itemize}
    \item \textbf{Contractions:} ``don't'' $\to$ ``do'' + ``n't'' or ``don't''?
    \item \textbf{Hyphenated words:} ``state-of-the-art'' --- one token or four?
    \item \textbf{URLs and emails:} \texttt{https://example.com} contains colons, slashes,
          and dots that na\"ive splitting would break on.
    \item \textbf{Numbers and punctuation:} ``\$3.14'' and ``U.S.A.'' require special
          handling.
    \item \textbf{CJK languages:} Chinese, Japanese, and Korean do not use whitespace
          between words. Chinese word segmentation is itself a major NLP task.
    \item \textbf{Agglutinative languages:} In Turkish or Finnish, a single whitespace-
          delimited token can encode an entire English sentence
          (``\emph{evlerinizden}'' = ``from your houses'').
\end{itemize}

Rule-based tokenizers (e.g., NLTK's \texttt{TreebankWordTokenizer}, spaCy's tokenizer)
use language-specific rules, regular expressions, and exception lists. They handle
common cases but require maintenance and fail on novel patterns.

\subsubsection{Practical Tokenization Challenges}

\begin{keyinsight}[The Tokenization Principle]
The choice of tokenizer has downstream consequences on vocabulary size, model capacity
allocation, and even what the model can learn. A tokenizer that splits ``unhappiness'' into
``un'' + ``happiness'' preserves morphological structure. One that splits it into ``unhap''
+ ``piness'' destroys it. This is not a trivial engineering detail---it is a modeling
decision.
\end{keyinsight}

The challenges listed above motivate the subword tokenization methods (BPE, WordPiece,
Unigram/SentencePiece) covered extensively in Chapter~3. In this chapter, we focus on the
classical word-level and character-level approaches that dominated the statistical NLP era.

\subsection{Encoding and Unicode Normalization}
\label{sec:unicode}

Real-world text comes in many encodings (UTF-8, UTF-16, Latin-1, etc.) and Unicode
normalization forms. Failure to handle encoding correctly leads to ``mojibake''---garbled
text that silently degrades model performance.

\begin{itemize}
    \item \textbf{Encoding normalization:} Convert all text to UTF-8. Use libraries like
    \texttt{chardet} or \texttt{cchardet} for detection, and \texttt{ftfy} (``fixes text
    for you'') for repair.

    \item \textbf{Unicode normalization (NFKC):} Unicode allows multiple representations
    of the same character. For example, the accented character ``\'{e}'' can be encoded as
    a single code point (\texttt{U+00E9}) or as ``e'' + combining accent
    (\texttt{U+0065 U+0301}). NFKC normalization maps both to a canonical form and also
    handles compatibility equivalences (e.g., ligatures, width variants).

    \item \textbf{Practical advice:} Always normalize to NFKC before any text processing.
    In Python: \texttt{unicodedata.normalize('NFKC', text)}.
\end{itemize}

\subsection{Case Normalization}
\label{sec:case_normalization}

Lowercasing reduces vocabulary size and collapses semantically identical terms (``The'' =
``the'' = ``THE''). However, case carries information:

\begin{itemize}
    \item \textbf{Named Entity Recognition:} Capitalization is a strong signal.
    ``Apple'' (company) vs. ``apple'' (fruit) is disambiguated partly by case.
    \item \textbf{Proper nouns:} ``polish'' (verb) vs. ``Polish'' (nationality).
    \item \textbf{Abbreviations and acronyms:} ``US'' vs. ``us''.
\end{itemize}

\begin{warningbox}
Blindly lowercasing all text is one of the most common preprocessing mistakes. For
information retrieval and topic classification, lowercasing is usually beneficial. For
NER, sentiment analysis of product names, or any task where capitalization is informative,
it can destroy critical signal. The correct answer in an interview is: ``it depends on the
task, and I would measure the impact empirically.''
\end{warningbox}

\subsection{Stopword Removal}
\label{sec:stopword_removal}

Stopwords are high-frequency, low-content words (``the,'' ``is,'' ``at,'' ``of''). Removing
them reduces dimensionality for bag-of-words models and speeds up retrieval.

\textbf{Arguments for removal:}
\begin{itemize}
    \item Reduces feature space dimensionality for classical models.
    \item Focuses TF-IDF and BoW on content-bearing words.
    \item Speeds up inverted index lookups in search systems.
\end{itemize}

\textbf{Arguments against removal:}
\begin{itemize}
    \item Destroys meaning in phrases: ``to be or not to be'' becomes empty after stopword
    removal.
    \item Harms performance on tasks sensitive to function words (authorship attribution,
    paraphrase detection).
    \item Modern neural models and subword tokenizers handle high-frequency words naturally;
    stopword removal is unnecessary and sometimes harmful.
    \item Negation words (``not,'' ``no,'' ``never'') are often on stopword lists but carry
    crucial sentiment information.
\end{itemize}

\begin{keyinsight}[When to Remove Stopwords]
Use stopword removal with classical bag-of-words and TF-IDF models for tasks like document
retrieval and topic classification. Do \emph{not} remove stopwords when using neural
models, when word order matters, or when function words carry task-relevant signal.
\end{keyinsight}

\subsection{Stemming vs. Lemmatization}
\label{sec:stemming_lemma}

Both stemming and lemmatization reduce words to a base form to consolidate vocabulary, but
they differ fundamentally in approach and quality.

\subsubsection{Stemming}

Stemming applies heuristic rules to strip affixes. It is fast, requires no dictionary, but
produces non-words.

\begin{itemize}
    \item \textbf{Porter Stemmer (1980):} The most widely used English stemmer. Applies a
    cascade of ~60 rules in five phases. Example: ``running'' $\to$ ``run'',
    ``generalization'' $\to$ ``gener'', ``studies'' $\to$ ``studi''.
    \item \textbf{Snowball Stemmer:} Improved, multilingual version of Porter. Available for
    ~25 languages.
    \item \textbf{Lancaster Stemmer:} More aggressive than Porter. ``maximum'' $\to$
    ``maxim'', ``presumably'' $\to$ ``presum''.
\end{itemize}

\subsubsection{Lemmatization}

Lemmatization uses vocabulary lookup and morphological analysis to reduce words to their
dictionary form (\emph{lemma}).

\begin{itemize}
    \item \textbf{WordNet Lemmatizer:} Looks up words in the WordNet lexical database.
    Requires POS tag for accuracy. ``better'' $\to$ ``good'' (adj), ``ran'' $\to$
    ``run'' (verb).
    \item \textbf{spaCy Lemmatizer:} Uses rule tables + exceptions. Fast and accurate for
    supported languages.
\end{itemize}

\subsubsection{Comparison}

\begin{table}[H]
\centering
\caption{Stemming vs. lemmatization: key differences.}
\label{tab:stem_vs_lemma}
\begin{tabular}{@{}lll@{}}
\toprule
\textbf{Aspect} & \textbf{Stemming} & \textbf{Lemmatization} \\
\midrule
Approach & Heuristic rules (affix stripping) & Dictionary + morphological analysis \\
Output & Often non-words (``studi'', ``gener'') & Valid dictionary words (``study'', ``general'') \\
Speed & Very fast (no lookup) & Slower (requires dictionary/POS) \\
Accuracy & Approximate & Precise \\
Languages & Easy to extend (rule-based) & Requires language-specific resources \\
Requires POS? & No & Often yes (for disambiguation) \\
Use cases & IR, search index, fast prototyping & Text analytics, NLU, display \\
\bottomrule
\end{tabular}
\end{table}

\subsection{Sentence Segmentation}
\label{sec:sentence_seg}

Splitting text into sentences seems trivial until you encounter real data:

\begin{itemize}
    \item Abbreviations: ``Dr. Smith went to Washington D.C. on Jan. 5th.'' contains four
    periods, only one of which ends a sentence.
    \item Ellipses: ``Wait\ldots{} what?'' --- how many sentences?
    \item Numbers: ``The price is \$3.50.'' --- the period after ``50'' is a sentence
    boundary, but the one after ``3'' is not.
    \item Quoted speech: ```He said, ``Go away.'' She left.' '' --- nested quotation with
    sentence boundaries inside.
\end{itemize}

Practical solutions include rule-based segmenters (NLTK's Punkt tokenizer, which uses
unsupervised abbreviation detection), spaCy's dependency-parse-based segmenter, and
regular expression heuristics with exception lists.

% ============================================================================
\section{Classical Text Representations}
\label{sec:representations}
% ============================================================================

Before the neural revolution, NLP required explicit methods to convert text into numerical
vectors. These classical representations remain foundational: they are still used in
production, they illuminate the principles that neural methods build upon, and they are
directly tested in interviews.

\subsection{One-Hot Encoding}
\label{sec:one_hot}

The simplest word representation assigns each word in the vocabulary a unique index and
represents it as a binary vector with a single 1.

For a vocabulary of size $V$, word $w_i$ is represented as:
\begin{equation}
    \mathbf{e}_i \in \{0, 1\}^V, \quad [\mathbf{e}_i]_j = \begin{cases} 1 & \text{if } j = i \\ 0 & \text{otherwise} \end{cases}
\end{equation}

\textbf{Limitations:}
\begin{itemize}
    \item \textbf{Extreme sparsity:} Vectors are $V$-dimensional with a single non-zero
    entry. For $V = 100{,}000$, this is 99.999\% zeros.
    \item \textbf{No similarity:} $\inner{\mathbf{e}_i}{\mathbf{e}_j} = 0$ for all $i
    \neq j$, so ``cat'' and ``dog'' are as different as ``cat'' and ``quantum.''
    \item \textbf{Vocabulary dependence:} Adding a new word requires changing the
    dimensionality of every vector.
\end{itemize}

One-hot encoding serves as a lookup index into embedding matrices rather than as a
meaningful representation in its own right.

\subsection{Bag of Words (BoW)}
\label{sec:bow}

The Bag of Words model represents a document as a vector of word counts, discarding word
order entirely.

Given a vocabulary $\mathcal{V} = \{w_1, w_2, \ldots, w_V\}$ and a document $d$, the BoW
representation is:
\begin{equation}
    \vx_d = \bigl[\text{count}(w_1, d), \; \text{count}(w_2, d), \; \ldots, \; \text{count}(w_V, d)\bigr] \in \N^V
\end{equation}

\textbf{Properties:}
\begin{itemize}
    \item \textbf{Simple and fast:} Requires only counting, no learned parameters.
    \item \textbf{Sparse:} Most entries are zero (a typical document uses a tiny fraction
    of the vocabulary).
    \item \textbf{Loses word order:} ``Dog bites man'' and ``Man bites dog'' have identical
    BoW representations.
    \item \textbf{No semantics:} Like one-hot, it captures no similarity between words.
    \item \textbf{High-frequency bias:} Common words dominate the representation.
\end{itemize}

Despite these limitations, BoW combined with the right classifier (Naive Bayes, SVM, or
logistic regression) remains a strong baseline for text classification, especially when
training data is limited or latency requirements are tight.

\subsection{N-Grams}
\label{sec:ngrams}

N-grams extend the BoW by capturing local word order. An $n$-gram is a contiguous
sequence of $n$ tokens from a text.

\begin{definition}[N-Gram]
Given a sequence of tokens $(t_1, t_2, \ldots, t_m)$, the set of $n$-grams is:
\begin{equation}
    \text{N-grams}(n) = \{(t_i, t_{i+1}, \ldots, t_{i+n-1}) \mid i = 1, \ldots, m - n + 1\}
\end{equation}
\end{definition}

\textbf{Examples for ``the cat sat on the mat'':}
\begin{itemize}
    \item \textbf{Unigrams ($n=1$):} the, cat, sat, on, the, mat
    \item \textbf{Bigrams ($n=2$):} the cat, cat sat, sat on, on the, the mat
    \item \textbf{Trigrams ($n=3$):} the cat sat, cat sat on, sat on the, on the mat
\end{itemize}

\textbf{Key tradeoff:} Larger $n$ captures more context but suffers from exponential
growth. For a vocabulary of size $V$, the number of possible $n$-grams is $V^n$. With
$V = 100{,}000$:
\begin{itemize}
    \item Unigrams: $10^5$ possible features
    \item Bigrams: $10^{10}$ possible features
    \item Trigrams: $10^{15}$ possible features
\end{itemize}

In practice, $n \leq 3$ is typical, with aggressive pruning of rare $n$-grams. The data
sparsity problem is fundamental: most possible $n$-grams never appear in any corpus.

\begin{keyinsight}[N-Grams in the Age of Transformers]
N-gram features are still used in (1) BM25 retrieval systems as a fast pre-filtering step,
(2) language identification models, (3) spelling correction, (4) text classification
pipelines where latency is critical, and (5) as evaluation metrics (BLEU is based on
$n$-gram precision). Interviewers value candidates who know when simpler methods suffice.
\end{keyinsight}

\subsection{TF-IDF: Term Frequency--Inverse Document Frequency}
\label{sec:tfidf}

TF-IDF is the most important classical text representation for information retrieval.
It addresses a fundamental flaw in raw term frequency: common words (``the,'' ``is,''
``and'') have high counts in every document but carry no discriminative information.

\subsubsection{Derivation}

\textbf{Term Frequency (TF):} The count of term $t$ in document $d$:
\begin{equation}
    \text{tf}(t, d) = f_{t,d}
    \label{eq:tf_raw}
\end{equation}
where $f_{t,d}$ is the raw count of term $t$ in document $d$.

\textbf{Problem with raw TF:} A document mentioning ``machine'' 10 times is not 10 times
more relevant than one mentioning it once. This motivates \textbf{sublinear TF}:
\begin{equation}
    \text{tf}_{\log}(t, d) = \begin{cases} 1 + \log f_{t,d} & \text{if } f_{t,d} > 0 \\ 0 & \text{otherwise} \end{cases}
    \label{eq:tf_sublinear}
\end{equation}

\textbf{Inverse Document Frequency (IDF):} Measures how informative a term is across
the corpus. A term that appears in every document has zero discriminative value.
\begin{equation}
    \text{idf}(t) = \log \frac{N}{\text{df}(t)}
    \label{eq:idf}
\end{equation}
where $N$ is the total number of documents and $\text{df}(t)$ is the number of documents
containing term $t$.

\textbf{Smoothed IDF variant} (avoids division by zero and reduces the impact of very
rare terms):
\begin{equation}
    \text{idf}_{\text{smooth}}(t) = \log \frac{N + 1}{\text{df}(t) + 1} + 1
    \label{eq:idf_smooth}
\end{equation}

\begin{mathresult}{TF-IDF Formula}
The TF-IDF weight of term $t$ in document $d$ from a corpus of $N$ documents is:
\begin{equation}
    \text{tfidf}(t, d) = \text{tf}(t, d) \times \text{idf}(t) = f_{t,d} \cdot \log \frac{N}{\text{df}(t)}
    \label{eq:tfidf}
\end{equation}

With sublinear TF scaling:
\begin{equation}
    \text{tfidf}_{\log}(t, d) = (1 + \log f_{t,d}) \cdot \log \frac{N}{\text{df}(t)}
    \label{eq:tfidf_sublinear}
\end{equation}

A document is represented as a vector $\vx_d \in \R^V$ where each dimension is the
TF-IDF weight of the corresponding vocabulary term. Similarity between documents is
typically computed using cosine similarity:
\begin{equation}
    \text{sim}(d_1, d_2) = \frac{\vx_{d_1} \cdot \vx_{d_2}}{\norm{\vx_{d_1}} \cdot \norm{\vx_{d_2}}}
\end{equation}
\end{mathresult}

\subsubsection{Why TF-IDF Works}

The key insight is the interaction between the two components:

\begin{itemize}
    \item \textbf{TF upweights} terms that appear frequently in a specific document
    (local signal).
    \item \textbf{IDF downweights} terms that appear frequently across the corpus
    (global signal).
    \item The product identifies terms that are \emph{locally frequent but globally rare}
    ---exactly the terms that discriminate one document from others.
\end{itemize}

A word like ``the'' has high TF in every document but near-zero IDF ($\text{df}
\approx N$). A word like ``transformer'' has moderate TF in relevant documents and high IDF
(appears in few documents). TF-IDF surfaces ``transformer'' as the discriminative term.

\subsubsection{Limitations of TF-IDF}

\begin{itemize}
    \item \textbf{No semantic similarity:} ``car'' and ``automobile'' have zero similarity
    in TF-IDF space because they are different vocabulary entries.
    \item \textbf{No word order:} Like BoW, word order is lost entirely.
    \item \textbf{Vocabulary mismatch:} A query about ``ML'' will not match a document
    about ``machine learning'' unless both terms appear.
    \item \textbf{Sparse and high-dimensional:} Vectors live in $\R^V$ with most
    entries zero.
\end{itemize}

\subsection{BM25: Best Matching 25}
\label{sec:bm25}

BM25 (Robertson et al., 1994) is the \emph{de facto} standard for lexical retrieval in
production search systems. It improves upon TF-IDF with two critical modifications:
term frequency saturation and document length normalization.

\subsubsection{Motivation}

TF-IDF has two problems for ranking:
\begin{enumerate}
    \item \textbf{No TF saturation:} Doubling the term count doubles the score, but
    relevance does not increase linearly with count. The 100th mention of ``cat'' in a
    document adds almost no additional evidence of relevance.
    \item \textbf{No length normalization:} Longer documents naturally contain more term
    occurrences and receive higher scores, regardless of actual relevance.
\end{enumerate}

BM25 addresses both issues.

\subsubsection{The BM25 Formula}

\begin{mathresult}{BM25 Scoring Function}
For a query $Q = \{q_1, q_2, \ldots, q_m\}$ and a document $d$:
\begin{equation}
    \text{BM25}(Q, d) = \sum_{i=1}^{m} \text{idf}(q_i) \cdot \frac{f_{q_i, d} \cdot (k_1 + 1)}{f_{q_i, d} + k_1 \cdot \left(1 - b + b \cdot \dfrac{|d|}{d_{\text{avg}}}\right)}
    \label{eq:bm25}
\end{equation}
where:
\begin{itemize}[itemsep=4pt]
    \item $f_{q_i, d}$ = frequency of query term $q_i$ in document $d$
    \item $|d|$ = document length (in tokens)
    \item $d_{\text{avg}}$ = average document length across the corpus
    \item $k_1$ = term frequency saturation parameter (typically $k_1 \in [1.2, 2.0]$)
    \item $b$ = length normalization parameter (typically $b = 0.75$)
    \item $\text{idf}(q_i)$ = inverse document frequency, commonly:
\end{itemize}
\begin{equation}
    \text{idf}(q_i) = \log \frac{N - \text{df}(q_i) + 0.5}{\text{df}(q_i) + 0.5}
    \label{eq:bm25_idf}
\end{equation}
\end{mathresult}

\subsubsection{Understanding the BM25 Parameters}

\textbf{Term Frequency Saturation ($k_1$):}

The BM25 TF component has the form:
\begin{equation}
    \text{tf}_{\text{BM25}} = \frac{f \cdot (k_1 + 1)}{f + k_1}
    \label{eq:bm25_tf}
\end{equation}

As $f \to \infty$, $\text{tf}_{\text{BM25}} \to (k_1 + 1)$, creating a saturation
effect. Additional occurrences yield diminishing returns. This contrasts with raw TF,
which grows without bound.

\begin{itemize}
    \item When $k_1 = 0$: TF is ignored entirely (binary model---term present or not).
    \item When $k_1 \to \infty$: No saturation, approaches raw TF scaling.
    \item Typical $k_1 = 1.2$: Moderate saturation, works well empirically.
\end{itemize}

\textbf{Length Normalization ($b$):}

The denominator includes $k_1 \cdot (1 - b + b \cdot |d| / d_{\text{avg}})$, which
adjusts for document length:

\begin{itemize}
    \item When $b = 0$: No length normalization. Long documents are treated like short
    ones.
    \item When $b = 1$: Full length normalization. A document twice the average length
    requires twice the term frequency to achieve the same score.
    \item Typical $b = 0.75$: Partial normalization, balancing length bias and the
    legitimate higher term counts in longer (but more comprehensive) documents.
\end{itemize}

\begin{keyinsight}[Why BM25 Still Dominates Lexical Retrieval]
BM25 remains the backbone of Elasticsearch, Apache Solr, and most production search
systems. It is fast (inverted index lookup), interpretable (you can explain why a document
ranked highly), requires no training data, and is competitive with or complementary to
neural retrievers. The best modern retrieval systems are \emph{hybrid}: BM25 for lexical
matching plus dense retrieval for semantic matching. Knowing BM25 in detail signals
production experience to interviewers.
\end{keyinsight}

\subsection{Comparison of Classical Representations}
\label{sec:repr_comparison}

\begin{table}[H]
\centering
\caption{Comparison of classical text representations.}
\label{tab:classical_repr}
\small
\begin{tabular}{@{}p{2.2cm}p{1.6cm}p{1.6cm}p{1.4cm}p{2.2cm}p{3.8cm}@{}}
\toprule
\textbf{Method} & \textbf{Word Order} & \textbf{Semantic Similarity} & \textbf{Sparse} & \textbf{Training Required} & \textbf{Best Use Case} \\
\midrule
One-Hot & None & None & Extreme & No & Embedding lookup index \\
BoW & None & None & High & No & Simple classification \\
N-Grams & Local ($n$) & None & Very High & No & LangID, spelling, BLEU \\
TF-IDF & None & None & High & No (corpus stats) & Doc similarity, classification \\
BM25 & None & None & High & No (corpus stats) & Production search ranking \\
\bottomrule
\end{tabular}
\end{table}

% ============================================================================
\section{Syntactic Analysis Foundations}
\label{sec:syntax}
% ============================================================================

Syntactic analysis concerns the grammatical structure of sentences. While modern LLMs
implicitly capture syntactic patterns, explicit syntactic analysis remains important for
structured information extraction, semantic parsing, and understanding what neural models
learn. Interviewers test syntax knowledge to assess foundational CS and linguistics
understanding.

\subsection{Part-of-Speech Tagging}
\label{sec:pos_tagging}

Part-of-Speech (POS) tagging assigns a grammatical category to each word in a sentence.

\subsubsection{The Penn Treebank Tagset}

The Penn Treebank tagset defines 36 POS tags for English, including:

\begin{table}[H]
\centering
\caption{Selected Penn Treebank POS tags.}
\label{tab:pos_tags}
\small
\begin{tabular}{@{}llll@{}}
\toprule
\textbf{Tag} & \textbf{Description} & \textbf{Tag} & \textbf{Description} \\
\midrule
NN & Noun, singular & VB & Verb, base form \\
NNS & Noun, plural & VBD & Verb, past tense \\
NNP & Proper noun, singular & VBG & Verb, gerund \\
JJ & Adjective & RB & Adverb \\
DT & Determiner & IN & Preposition \\
PRP & Personal pronoun & CC & Coordinating conjunction \\
\bottomrule
\end{tabular}
\end{table}

\subsubsection{Why POS Tagging Matters}

\begin{itemize}
    \item \textbf{Disambiguates word usage:} ``book'' can be a noun (NN) or verb (VB);
    POS reveals the intended usage.
    \item \textbf{Improves downstream tasks:} NER accuracy increases with POS features
    (proper nouns are often entities). Lemmatization requires POS (``better'' $\to$
    ``good'' only as an adjective).
    \item \textbf{Feature engineering:} POS patterns like ``JJ NN'' (adjective + noun)
    identify noun phrases for information extraction.
\end{itemize}

\subsubsection{Approaches}

The evolution of POS tagging mirrors the broader NLP paradigm shifts:
\begin{enumerate}
    \item \textbf{Rule-based:} Hand-crafted rules (e.g., if word ends in ``-ing,'' tag as
    VBG). Brittle, low coverage.
    \item \textbf{HMM (Hidden Markov Model):} Models $P(\text{tag}_i \given
    \text{tag}_{i-1}) \cdot P(\text{word}_i \given \text{tag}_i)$. Transition and
    emission probabilities learned from data. Decoding via Viterbi algorithm.
    \item \textbf{CRF (Conditional Random Field):} Discriminative model that directly
    models $P(\text{tags} \given \text{words})$. Richer features than HMM.
    \item \textbf{BiLSTM-CRF:} Neural features from bidirectional LSTM fed into a CRF
    layer. Strong baseline through 2018.
    \item \textbf{Transformer-based:} BERT + linear classification layer. Current
    state-of-the-art (~97.5\% accuracy on Penn Treebank).
\end{enumerate}

\subsection{Named Entity Recognition (NER)}
\label{sec:ner}

Named Entity Recognition identifies and classifies named entities in text into predefined
categories.

\subsubsection{Standard Entity Types}

\begin{itemize}
    \item \textbf{PER:} Person names (``Barack Obama'')
    \item \textbf{ORG:} Organizations (``Google,'' ``United Nations'')
    \item \textbf{LOC:} Locations (``Paris,'' ``Mount Everest'')
    \item \textbf{MISC:} Miscellaneous entities (``English,'' ``World Cup'')
    \item Domain-specific types: DATE, TIME, MONEY, PRODUCT, EVENT, etc.
\end{itemize}

\subsubsection{BIO and BIOES Tagging Schemes}

NER is formulated as a sequence labeling problem using special tagging schemes that
encode entity boundaries.

\begin{figure}[H]
\centering
\begin{tikzpicture}[
    word/.style={font=\small\ttfamily, minimum width=1.2cm, minimum height=0.5cm},
    tag/.style={font=\small\sffamily\bfseries, minimum width=1.2cm, minimum height=0.5cm}
]
    % Sentence
    \node[font=\small\sffamily\bfseries, anchor=west] at (-2.5, 1.2) {Sentence:};
    \node[font=\small\sffamily\bfseries, anchor=west] at (-2.5, 0.4) {BIO Tags:};
    \node[font=\small\sffamily\bfseries, anchor=west] at (-2.5, -0.4) {BIOES Tags:};

    % Words
    \foreach \w/\i in {Barack/0, Obama/1, visited/2, New/3, York/4, City/5, yesterday/6}{
        \node[word] (w\i) at (1.3*\i + 0.5, 1.2) {\w};
    }
    % BIO tags
    \foreach \t/\c/\i in {B-PER/deepblue/0, I-PER/deepblue/1, O/gray/2, B-LOC/deepgreen/3, I-LOC/deepgreen/4, I-LOC/deepgreen/5, O/gray/6}{
        \node[tag, text=\c] at (1.3*\i + 0.5, 0.4) {\t};
    }
    % BIOES tags
    \foreach \t/\c/\i in {B-PER/deepblue/0, E-PER/deepblue/1, O/gray/2, B-LOC/deepgreen/3, I-LOC/deepgreen/4, E-LOC/deepgreen/5, O/gray/6}{
        \node[tag, text=\c] at (1.3*\i + 0.5, -0.4) {\t};
    }

    % Entity spans
    \draw[thick, deepblue, decorate, decoration={brace, amplitude=4pt, mirror}]
        (w0.south west) ++(0,-1.1) -- (w1.south east) ++(0,-1.1)
        node[midway, below=5pt, font=\scriptsize\sffamily, text=deepblue] {PER};
    \draw[thick, deepgreen, decorate, decoration={brace, amplitude=4pt, mirror}]
        (w3.south west) ++(0,-1.1) -- (w5.south east) ++(0,-1.1)
        node[midway, below=5pt, font=\scriptsize\sffamily, text=deepgreen] {LOC};
\end{tikzpicture}
\caption{BIO and BIOES tagging schemes for NER. B = Beginning, I = Inside, O = Outside,
E = End, S = Single (singleton entity, not shown here).}
\label{fig:bio_tagging}
\end{figure}

\textbf{BIO scheme:}
\begin{itemize}
    \item \textbf{B-X:} Beginning of entity of type X.
    \item \textbf{I-X:} Inside (continuation of) entity of type X.
    \item \textbf{O:} Outside any entity.
\end{itemize}

\textbf{BIOES scheme} (also called BILOU) adds:
\begin{itemize}
    \item \textbf{E-X:} End of entity of type X (last token of a multi-token entity).
    \item \textbf{S-X:} Singleton entity of type X (single-token entity).
\end{itemize}

\begin{keyinsight}[Why BIOES Is Better Than BIO]
BIOES provides more explicit boundary information. In BIO, the model must infer entity
boundaries from transitions (B-PER $\to$ I-PER vs. B-PER $\to$ B-PER for adjacent
entities). BIOES makes boundaries explicit, which is especially helpful for (1) adjacent
entities of the same type (``New York'' [LOC] next to ``City Hall'' [LOC]) and (2)
single-token entities. Empirically, BIOES improves NER F1 by 0.5--1.5\% over BIO.
\end{keyinsight}

\subsubsection{Why Not Classify Each Token Independently?}

This is a critical question that reveals understanding of structured prediction.
Independent token classification with a softmax layer makes predictions for each position
without considering the labels of neighboring tokens. This allows illegal sequences:
\begin{itemize}
    \item \texttt{I-PER} following \texttt{B-LOC} (entity type switch without a beginning
    tag).
    \item \texttt{I-PER} at the start of a sequence (inside tag without a beginning).
    \item Inconsistent entity boundaries.
\end{itemize}

A CRF layer on top of the neural encoder models the \emph{joint} probability of the
entire label sequence, enforcing transition constraints (e.g., \texttt{I-PER} can only
follow \texttt{B-PER} or \texttt{I-PER}). The Viterbi algorithm then finds the globally
optimal label sequence rather than a sequence of locally optimal labels.

\subsection{Constituency vs. Dependency Parsing}
\label{sec:parsing}

Parsing reveals the grammatical structure of sentences. Two dominant formalisms exist:

\subsubsection{Constituency Parsing}

Constituency parsing decomposes sentences into nested constituents (phrases) according to a
context-free grammar (CFG). Each internal node has a phrase label (S, NP, VP, PP, etc.).

\begin{figure}[H]
\centering
\begin{tikzpicture}[
    level distance=1.1cm,
    sibling distance=2cm,
    every node/.style={font=\small\sffamily},
    edge from parent/.style={draw, -},
    level 1/.style={sibling distance=4cm},
    level 2/.style={sibling distance=2cm}
]
\node {S}
    child { node {NP}
        child { node {DT}
            child { node {\textit{The}} }
        }
        child { node {NN}
            child { node {\textit{cat}} }
        }
    }
    child { node {VP}
        child { node {VBD}
            child { node {\textit{sat}} }
        }
        child { node {PP}
            child { node {IN}
                child { node {\textit{on}} }
            }
            child { node {NP}
                child { node {DT}
                    child { node {\textit{the}} }
                }
                child { node {NN}
                    child { node {\textit{mat}} }
                }
            }
        }
    };
\end{tikzpicture}
\caption{Constituency parse tree for ``The cat sat on the mat.'' Non-terminal nodes
represent phrase categories; terminal nodes are words.}
\label{fig:constituency_parse}
\end{figure}

\textbf{Key algorithms:} CKY (Cocke--Kasami--Younger) for exact parsing with CFGs,
neural parsers using pointer networks for more flexible parsing.

\subsubsection{Dependency Parsing}

Dependency parsing represents grammatical structure as directed edges (dependencies)
between words. Each word is either the root of the sentence or a dependent of exactly
one other word (its head). Labels indicate the grammatical relation (nsubj, dobj, prep,
etc.).

\begin{figure}[H]
\centering
\begin{tikzpicture}[
    word/.style={font=\small\ttfamily, anchor=base},
    dep/.style={-{Stealth[length=2mm]}, thick, deepblue},
    label/.style={font=\scriptsize\sffamily\itshape, text=deepgreen, fill=white, inner sep=1pt}
]
    % Words
    \node[word] (the1) at (0,0) {The};
    \node[word] (cat) at (1.5,0) {cat};
    \node[word] (sat) at (3,0) {sat};
    \node[word] (on) at (4.5,0) {on};
    \node[word] (the2) at (6,0) {the};
    \node[word] (mat) at (7.5,0) {mat};

    % Dependencies (arcs above)
    \draw[dep] (cat.north) to[out=130,in=50] node[label, above, pos=0.5] {det} (the1.north);
    \draw[dep] (sat.north) to[out=130,in=50] node[label, above, pos=0.5] {nsubj} (cat.north);
    \draw[dep] (sat.north) to[out=50,in=130] node[label, above, pos=0.5] {prep} (on.north);
    \draw[dep] (on.north) to[out=50,in=130] node[label, above, pos=0.5] {pobj} (mat.north);
    \draw[dep] (mat.north) to[out=130,in=50] node[label, above, pos=0.5] {det} (the2.north);

    % Root
    \draw[dep, deepred] (sat.north) to[out=90,in=-90] ++(0,1.2)
        node[above, font=\scriptsize\sffamily\bfseries, text=deepred] {ROOT};
\end{tikzpicture}
\caption{Dependency parse of ``The cat sat on the mat.'' Arcs point from heads to
dependents with grammatical relation labels.}
\label{fig:dependency_parse}
\end{figure}

\textbf{Key approaches:}
\begin{itemize}
    \item \textbf{Transition-based (Shift-Reduce):} Processes tokens left-to-right using a
    stack and a buffer. Actions: SHIFT (push token onto stack), LEFT-ARC (create
    dependency left), RIGHT-ARC (create dependency right). Fast: $O(n)$ time.
    \item \textbf{Graph-based (MST):} Scores all possible edges, finds the maximum
    spanning tree. Uses Eisner's algorithm for projective parsing ($O(n^3)$) or
    Chu--Liu--Edmonds for non-projective ($O(n^2)$). More accurate, slower.
\end{itemize}

\subsubsection{Constituency vs. Dependency: Comparison}

\begin{table}[H]
\centering
\caption{Constituency vs. dependency parsing comparison.}
\label{tab:parsing_comparison}
\small
\begin{tabular}{@{}lll@{}}
\toprule
\textbf{Aspect} & \textbf{Constituency} & \textbf{Dependency} \\
\midrule
Structure & Nested phrases (tree) & Head-dependent relations (graph) \\
Labels & Phrase categories (NP, VP) & Grammatical relations (nsubj, dobj) \\
Nodes & Words + phrase nodes & Words only \\
Best for & Phrase structure, CFG applications & Relation extraction, semantic roles \\
Free word order & Struggles & Handles well \\
Complexity & $O(n^3)$ with CKY & $O(n)$ transition / $O(n^3)$ graph \\
\bottomrule
\end{tabular}
\end{table}

\subsection{Why Syntax Still Matters}
\label{sec:why_syntax}

Even in the LLM era, syntactic knowledge remains relevant:

\begin{enumerate}
    \item \textbf{Information Extraction:} Dependency paths between entities are strong
    features for relation extraction (``X, CEO of Y'' follows a consistent syntactic
    pattern).
    \item \textbf{Semantic Parsing:} Converting natural language to structured queries
    (SQL, SPARQL) often uses syntactic structure as an intermediate representation.
    \item \textbf{Understanding Attention Patterns:} Research shows that certain attention
    heads in BERT and GPT learn syntactic patterns resembling dependency relations
    (Clark et al., 2019). Understanding syntax helps interpret what transformers learn.
    \item \textbf{Grammar Checking:} Rule-based and hybrid grammar checkers (Grammarly's
    early systems, LanguageTool) rely on syntactic parsing.
    \item \textbf{Low-resource and Specialized Domains:} When LLMs are unavailable
    (on-device, privacy constraints, domain-specific), syntactic parsing provides
    structured analysis with smaller models.
\end{enumerate}

% ============================================================================
\section{N-Gram Language Models}
\label{sec:ngram_lm}
% ============================================================================

Language models assign probabilities to sequences of words. N-gram language models are the
classical approach, and understanding them is essential because they illuminate fundamental
concepts (chain rule, Markov assumption, smoothing, perplexity) that remain central to
modern NLP.

\subsection{The Chain Rule of Probability}
\label{sec:chain_rule}

The probability of a sequence of words $w_1, w_2, \ldots, w_T$ decomposes via the chain
rule:
\begin{equation}
    P(w_1, w_2, \ldots, w_T) = \prod_{t=1}^{T} P(w_t \given w_1, \ldots, w_{t-1})
    \label{eq:chain_rule}
\end{equation}

This is exact but intractable: estimating $P(w_t \given w_1, \ldots, w_{t-1})$ requires
conditioning on the \emph{entire} preceding context, which creates a combinatorial
explosion.

\subsection{The Markov Assumption}
\label{sec:markov}

The $n$-gram approximation applies the Markov assumption: the probability of the next
word depends only on the previous $n - 1$ words:

\begin{mathresult}{N-Gram Language Model}
\begin{equation}
    P(w_t \given w_1, \ldots, w_{t-1}) \approx P(w_t \given w_{t-n+1}, \ldots, w_{t-1})
    \label{eq:ngram_lm}
\end{equation}

Estimated from counts via Maximum Likelihood Estimation (MLE):
\begin{equation}
    P_{\text{MLE}}(w_t \given w_{t-n+1}, \ldots, w_{t-1}) = \frac{\text{count}(w_{t-n+1}, \ldots, w_t)}{\text{count}(w_{t-n+1}, \ldots, w_{t-1})}
    \label{eq:ngram_mle}
\end{equation}
\end{mathresult}

\begin{itemize}
    \item \textbf{Unigram ($n=1$):} $P(w_t) = \text{count}(w_t) / N_{\text{total}}$.
    No context at all.
    \item \textbf{Bigram ($n=2$):} $P(w_t \given w_{t-1})$. One word of context.
    \item \textbf{Trigram ($n=3$):} $P(w_t \given w_{t-2}, w_{t-1})$. Two words of context.
\end{itemize}

\subsection{Smoothing}
\label{sec:smoothing}

The MLE estimate assigns zero probability to any $n$-gram not seen in the training corpus.
This is catastrophic: a single unseen $n$-gram makes the entire sentence probability zero.
Smoothing techniques redistribute probability mass to unseen events.

\begin{itemize}
    \item \textbf{Laplace (Add-1) Smoothing:} Add 1 to all counts. Simple but assigns too
    much mass to unseen events, especially with large vocabularies.
    \begin{equation}
        P_{\text{Laplace}}(w_t \given w_{t-1}) = \frac{\text{count}(w_{t-1}, w_t) + 1}{\text{count}(w_{t-1}) + V}
    \end{equation}

    \item \textbf{Add-$k$ Smoothing:} Generalization with tunable $k < 1$. Better but
    still ad hoc.

    \item \textbf{Interpolation:} Combine $n$-gram models of different orders:
    \begin{equation}
        P_{\text{interp}}(w_t \given w_{t-2}, w_{t-1}) = \lambda_3 P_3 + \lambda_2 P_2 + \lambda_1 P_1, \quad \sum_i \lambda_i = 1
    \end{equation}
    where $\lambda_i$ are learned on held-out data.

    \item \textbf{Kneser-Ney Smoothing:} The gold standard. Uses absolute discounting
    (subtract a fixed $d$ from each count) and a novel backoff distribution based on the
    \emph{number of distinct contexts} a word appears in (continuation probability), not
    raw frequency. ``Francisco'' has high unigram count but low continuation count (almost
    always follows ``San''), so Kneser-Ney correctly gives it low backoff probability.
\end{itemize}

\subsection{Perplexity}
\label{sec:perplexity}

Perplexity is the standard evaluation metric for language models, measuring how ``surprised''
the model is by test data.

\begin{mathresult}{Perplexity}
For a test sequence $w_1, w_2, \ldots, w_T$:
\begin{equation}
    \text{PP}(W) = P(w_1, w_2, \ldots, w_T)^{-1/T} = \exp\!\left(-\frac{1}{T} \sum_{t=1}^{T} \log P(w_t \given w_{<t})\right)
    \label{eq:perplexity}
\end{equation}

Equivalently, $\text{PP} = \exp(H)$ where $H$ is the cross-entropy of the model on the
test data. Lower perplexity indicates a better model.
\end{mathresult}

\textbf{Intuition:} Perplexity represents the effective number of equally likely choices
the model faces at each step---the weighted average branching factor. A perplexity of 100
means the model is as uncertain as if it were choosing uniformly among 100 options at
each token.

\begin{warningbox}
Perplexity is widely used but has important limitations: (1) it does not measure
fluency, coherence, or factual accuracy of generated text; (2) it is
vocabulary-dependent---models with different tokenizers have incomparable perplexities;
(3) it penalizes all errors equally, whether they are severe (wrong fact) or minor (synonym
choice). For LLM evaluation, downstream task performance and human evaluation are more
informative than perplexity alone.
\end{warningbox}

% ============================================================================
\section{Classical Text Classification}
\label{sec:classical_classification}
% ============================================================================

Before neural networks, text classification relied on hand-crafted features and classical
ML algorithms. These methods remain relevant as baselines, for resource-constrained
environments, and as interview topics that test fundamental ML understanding.

\subsection{Naive Bayes for Text}
\label{sec:naive_bayes}

Naive Bayes applies Bayes' theorem with a strong conditional independence assumption:

\begin{equation}
    P(c \given d) \propto P(c) \prod_{i=1}^{n} P(w_i \given c)
    \label{eq:naive_bayes}
\end{equation}

where $c$ is the class, $d$ is the document, and $w_i$ are the words in $d$. The
``naive'' assumption is that words are conditionally independent given the class.

\textbf{Multinomial Naive Bayes} models word counts with a multinomial distribution:
\begin{equation}
    P(w_i \given c) = \frac{\text{count}(w_i, c) + \alpha}{\sum_{w \in V} \text{count}(w, c) + \alpha V}
\end{equation}
where $\alpha$ is the Laplace smoothing parameter.

\textbf{Why it works despite the independence assumption:} The conditional independence
assumption is clearly wrong (words are not independent), but Naive Bayes is often
competitive because: (1) it only needs to get the ranking right (which class has higher
posterior), not the exact probabilities; (2) the misestimation errors tend to cancel out
across many features; (3) it is extremely sample-efficient and robust to irrelevant
features.

\subsection{SVMs for Text Classification}
\label{sec:svm_text}

Linear SVMs with TF-IDF features were the dominant text classification method for over a
decade.

\textbf{Why linear SVMs work so well for text:} Text data is high-dimensional and sparse.
In high-dimensional spaces, data is often linearly separable or near-separable. The SVM
finds the maximum-margin hyperplane, which generalizes well. Kernel tricks are rarely
needed---the linear kernel suffices and scales to millions of features.

\textbf{The classical NLP pipeline was:}
\begin{enumerate}
    \item Tokenize and preprocess text.
    \item Extract features (TF-IDF vectors, possibly with n-gram features).
    \item Train a linear SVM (or logistic regression).
    \item Evaluate with precision, recall, F1.
\end{enumerate}

\begin{keyinsight}[When Classical Methods Beat Neural]
Classical pipelines (TF-IDF + SVM/logistic regression) can outperform fine-tuned BERT when:
(1) training data is very small ($< 100$ examples); (2) latency requirements are extremely
tight ($< 1$ ms); (3) the task is simple (binary classification with clear keyword
signals); (4) interpretability is required (TF-IDF weights are directly inspectable). A
strong interview answer demonstrates knowing the right tool for the right situation, not
defaulting to the most complex model.
\end{keyinsight}

% ============================================================================
\section{Interview Questions}
\label{sec:classical_interview}
% ============================================================================

% --- Question 1: TF-IDF ---
\begin{interviewq}
\textbf{Q: What is TF-IDF? Derive the formula. When would you use it over neural embeddings?}

\badgeLfive{L5} \quad \badgetype{Conceptual} \quad \badgefreq{\freqhigh}

\stronganswer
TF-IDF stands for Term Frequency--Inverse Document Frequency. It is a weighting scheme that
assigns higher scores to terms that are frequent in a specific document but rare across the
corpus.

\textbf{Term Frequency} measures how often a term $t$ appears in document $d$:
$\text{tf}(t, d) = f_{t,d}$. A sublinear variant uses $1 + \log f_{t,d}$ to dampen the
effect of very high counts.

\textbf{Inverse Document Frequency} measures how informative a term is:
$\text{idf}(t) = \log(N / \text{df}(t))$, where $N$ is the total number of documents and
$\text{df}(t)$ is the number of documents containing $t$. Common words like ``the'' appear
in nearly every document, so their IDF approaches zero. Rare, discriminative terms have high
IDF.

The combined score is $\text{tfidf}(t, d) = \text{tf}(t, d) \times \text{idf}(t)$.

I would use TF-IDF over neural embeddings when: (1) I need interpretable features---each
dimension corresponds to a specific term; (2) latency is critical---TF-IDF is a simple
sparse vector operation with no GPU needed; (3) training data is limited---TF-IDF requires
no labeled data, only corpus statistics; (4) the task relies on exact lexical matching, such
as keyword-based search or duplicate detection where semantics are not the bottleneck.

\redflags
\begin{itemize}
    \item Saying ``TF-IDF is outdated and should never be used.''
    \item Not knowing the IDF formula or why the log is needed.
    \item Being unable to explain why common words get low scores.
    \item Forgetting to mention sublinear TF or document normalization.
\end{itemize}

\interviewtest{Whether the candidate understands classical IR fundamentals, can derive
formulas, and can reason about when simpler methods are appropriate.}

\goodgreat
\begin{itemize}
    \item \badgeLfive{L5}: Derives the formula correctly, explains the intuition, gives
    one use case for TF-IDF over embeddings.
    \item \badgeLsix{L6}: Discusses variants (sublinear TF, smoothed IDF), connects to BM25,
    explains the vocabulary mismatch limitation and how dense retrieval addresses it.
    \item \badgeLseven{L7}: Discusses hybrid retrieval (BM25 + dense), explains when
    TF-IDF features are used alongside neural models (e.g., as features in a learning-to-rank
    system), and how production search systems combine lexical and semantic signals.
\end{itemize}

\followups{How does BM25 improve on TF-IDF? What is the vocabulary mismatch problem?
How would you combine TF-IDF with neural methods in a hybrid retrieval system?}
\end{interviewq}

\vspace{0.5em}

% --- Question 2: BM25 ---
\begin{interviewq}
\textbf{Q: Explain BM25. Why is it still used in production search systems?}

\badgeLfive{L5} \quad \badgetype{Conceptual} \quad \badgefreq{\freqhigh}

\stronganswer
BM25 is a probabilistic retrieval function that scores the relevance of a document to a
query. It improves on TF-IDF in two key ways.

First, \textbf{term frequency saturation}: BM25 uses the formula $f(k_1+1)/(f+k_1)$ for
the TF component, which asymptotically approaches $k_1 + 1$ as term frequency $f$
increases. This means additional occurrences of a query term yield diminishing returns,
unlike raw TF which grows linearly.

Second, \textbf{document length normalization}: the denominator includes a factor
$k_1(1 - b + b \cdot |d|/d_{\text{avg}})$ that penalizes longer documents. A document
twice the average length needs roughly twice the term frequency to achieve the same score.
The parameter $b$ controls the strength of this normalization.

BM25 remains in production because it is: (1) \emph{fast}---implemented via inverted
index lookup, handling millions of documents in milliseconds; (2) \emph{requires no
training data}---only corpus statistics; (3) \emph{robust}---works well across domains
without tuning; (4) \emph{interpretable}---you can explain why a document ranked highly by
inspecting which query terms matched and their BM25 contributions; (5) \emph{complementary
to neural methods}---modern hybrid systems use BM25 for exact lexical matching alongside
dense retrieval for semantic matching.

\redflags
\begin{itemize}
    \item Not knowing the parameters $k_1$ and $b$ or their roles.
    \item Saying ``BM25 is just TF-IDF'' without explaining the differences.
    \item Not knowing that BM25 is still a production standard (Elasticsearch, Solr, Lucene).
\end{itemize}

\interviewtest{Depth of IR knowledge, understanding of why classical methods persist
in production, ability to reason about tradeoffs between simple and complex systems.}

\goodgreat
\begin{itemize}
    \item \badgeLfive{L5}: Explains the formula, both improvements over TF-IDF, and why BM25
    is fast.
    \item \badgeLsix{L6}: Discusses parameter tuning, the connection to probabilistic
    relevance models, and the BM25 IDF variant. Explains hybrid retrieval architecture.
    \item \badgeLseven{L7}: Discusses BM25F (field-weighted), learned sparse models like
    SPLADE that bridge BM25 and neural retrieval, and how modern search pipelines combine
    BM25 candidate retrieval with neural re-ranking.
\end{itemize}

\followups{What are typical values for $k_1$ and $b$? How does BM25 handle multi-word
queries? What is BM25F? How would you combine BM25 with a dense retriever?}
\end{interviewq}

\vspace{0.5em}

% --- Question 3: Stemming vs. Lemmatization ---
\begin{interviewq}
\textbf{Q: What is the difference between stemming and lemmatization? When does the choice matter?}

\badgeLfive{L5} \quad \badgetype{Conceptual} \quad \badgefreq{\freqmed}

\stronganswer
Both reduce words to a base form, but they differ in approach and output quality.

\textbf{Stemming} applies heuristic rules to strip suffixes. The Porter stemmer, for
example, uses about 60 rules in five phases. It is fast and requires no dictionary, but
produces non-words: ``studies'' becomes ``studi,'' ``generalization'' becomes ``gener.''
It over-stems (conflates unrelated words) and under-stems (fails to merge related words).

\textbf{Lemmatization} uses vocabulary lookup and morphological analysis to produce valid
dictionary forms: ``studies'' becomes ``study,'' ``better'' becomes ``good'' (as adjective),
``ran'' becomes ``run.'' It is more accurate but slower and typically requires POS
information for disambiguation.

The choice matters in several scenarios: (1) For \emph{search and IR}, stemming is usually
sufficient and faster---collapsing ``running/runs/ran'' to ``run'' improves recall. (2) For
\emph{text analytics and display}, lemmatization is preferred because results must be
human-readable. (3) For \emph{morphologically rich languages} like Finnish or Turkish, simple
suffix-stripping is inadequate; proper morphological analysis is needed. (4) For
\emph{neural models with subword tokenizers}, neither is necessary---BPE and WordPiece
handle morphological variation naturally.

\redflags
\begin{itemize}
    \item Confusing the two or thinking they produce identical outputs.
    \item Not knowing that lemmatization requires POS disambiguation.
    \item Saying ``always use lemmatization because it is better'' without considering speed
    or task requirements.
\end{itemize}

\interviewtest{Understanding of text normalization tradeoffs and practical experience
with NLP preprocessing.}

\goodgreat
\begin{itemize}
    \item \badgeLfive{L5}: Clear distinction, gives examples, knows when each is useful.
    \item \badgeLsix{L6}: Discusses specific stemmers (Porter, Snowball, Lancaster), mentions
    language-specific challenges, and explains why subword tokenization makes both less
    necessary.
    \item \badgeLseven{L7}: Discusses production considerations (stemming at index time
    vs. query time in search systems), interaction with BM25/TF-IDF, and multilingual
    stemming challenges.
\end{itemize}

\followups{What is the Porter stemmer's algorithm? Give an example where stemming hurts
performance. How do subword tokenizers relate to stemming and lemmatization?}
\end{interviewq}

\vspace{0.5em}

% --- Question 4: BIO Tagging ---
\begin{interviewq}
\textbf{Q: What is the BIO tagging scheme? Why not just classify each token independently?}

\badgeLfive{L5} \quad \badgetype{Conceptual} \quad \badgefreq{\freqmed}

\stronganswer
BIO is a sequence labeling scheme used in NER where each token is tagged as B (beginning of
an entity), I (inside/continuation of an entity), or O (outside any entity). Entity types
are appended, giving tags like B-PER, I-PER, O. For example, ``Barack Obama visited Paris''
would be tagged B-PER I-PER O B-LOC.

You cannot just classify each token independently because NER has \emph{structural
constraints}. Independent classification treats each position in isolation, which can
produce illegal sequences---for instance, I-PER following B-LOC (entity type switches
mid-span), or I-PER appearing at the start of a sequence without a preceding B-PER.

To enforce these constraints, we use a CRF (Conditional Random Field) layer on top of the
token encoder. The CRF models the joint probability of the entire label sequence, including
transition scores between adjacent labels. It learns that certain transitions are illegal
(e.g., I-PER after B-LOC should have a very negative transition score) and uses the Viterbi
algorithm to find the globally optimal sequence rather than a sequence of locally optimal
labels.

The BIOES variant adds E (end of entity) and S (singleton entity) tags, providing more
explicit boundary information and typically improving F1 by 0.5--1.5\%.

\redflags
\begin{itemize}
    \item Not knowing what B, I, and O stand for.
    \item Not understanding why independent classification produces invalid sequences.
    \item Being unable to explain how a CRF enforces structural constraints.
\end{itemize}

\interviewtest{Understanding of structured prediction, sequence labeling, and the
difference between independent and joint modeling.}

\goodgreat
\begin{itemize}
    \item \badgeLfive{L5}: Explains BIO clearly, gives the independent classification
    failure case, mentions CRFs.
    \item \badgeLsix{L6}: Explains the CRF loss function, Viterbi decoding, and the
    BiLSTM-CRF architecture. Compares BIO and BIOES with concrete examples.
    \item \badgeLseven{L7}: Discusses span-based NER as an alternative to sequence labeling,
    nested NER, and when CRF layers help on top of transformers (yes for low-data, marginal
    for large BERT models).
\end{itemize}

\followups{How does Viterbi decoding work? What is the time complexity? When would
you add a CRF on top of BERT? What is nested NER and how do you handle it?}
\end{interviewq}

\vspace{0.5em}

% --- Question 5: Rule-based vs. Statistical vs. Neural ---
\begin{interviewq}
\textbf{Q: Explain the difference between rule-based, statistical, and neural NLP. When would you still use rule-based methods?}

\badgeLfive{L5} \quad \badgetype{Conceptual} \quad \badgefreq{\freqhigh}

\stronganswer
\textbf{Rule-based NLP} uses hand-crafted patterns: regular expressions, grammars,
gazetteers, and decision trees written by domain experts. It is interpretable and
deterministic but brittle, expensive to maintain, and fails to generalize.

\textbf{Statistical NLP} learns patterns from data using probabilistic models: HMMs for
sequence labeling, n-gram language models, CRFs, TF-IDF for retrieval. Features are
hand-engineered but model parameters are learned. This era achieved broad coverage but
required significant feature engineering expertise.

\textbf{Neural NLP} learns both features and parameters end-to-end from data. From
Word2Vec and LSTMs to BERT and GPT, neural methods progressively eliminated manual feature
engineering and achieved state-of-the-art on virtually every NLP benchmark.

I would still use rule-based methods when: (1) \emph{Precision is critical and recall is
secondary}---e.g., extracting structured data from known formats (dates, phone numbers,
email addresses) where regex is 100\% precise. (2) \emph{The problem is well-defined and
closed}---e.g., parsing URLs, validating input formats, simple text normalization. (3)
\emph{Interpretability is required by regulation}---e.g., content moderation rules that must
be auditable. (4) \emph{Speed is paramount}---regex and rule engines are orders of
magnitude faster than neural models. (5) \emph{No training data exists}---rules can be
written from domain knowledge alone.

In practice, the best production systems are hybrid: rule-based preprocessing and
post-processing with neural models in between.

\redflags
\begin{itemize}
    \item Saying ``rule-based methods are obsolete'' without qualification.
    \item Not understanding the role of feature engineering in statistical NLP.
    \item Being unable to give a concrete scenario where rules are preferable.
\end{itemize}

\interviewtest{Breadth of NLP knowledge, practical judgment about when to use which
approach, and appreciation for the full toolbox.}

\goodgreat
\begin{itemize}
    \item \badgeLfive{L5}: Clear distinction between the three paradigms with examples and
    at least two use cases for rule-based methods.
    \item \badgeLsix{L6}: Discusses hybrid systems, gives production examples (e.g., PII
    redaction with regex + NER model), explains the maintenance cost of each approach.
    \item \badgeLseven{L7}: Discusses how to decide the boundary between rules and models in
    a production system, coverage vs. precision tradeoffs, and how LLMs are changing
    this landscape (e.g., replacing rule cascades with prompted models).
\end{itemize}

\followups{Give an example of a hybrid rule + ML system. How do you test and maintain
rules at scale? When does a rule-based system become too complex to maintain?}
\end{interviewq}

\vspace{0.5em}

% --- Question 6: Bag of Words Problems ---
\begin{interviewq}
\textbf{Q: What problems does the bag-of-words assumption create? How have subsequent methods addressed them?}

\badgeLfive{L5} \quad \badgetype{Conceptual} \quad \badgefreq{\freqmed}

\stronganswer
The bag-of-words assumption discards word order entirely, treating a document as an
unordered collection of word counts. This creates several problems:

\textbf{1. Loss of word order and compositionality:} ``Dog bites man'' and ``Man bites
dog'' have identical BoW representations despite opposite meanings. Negation is lost:
``not good'' and ``good'' differ by only a stopword.

\textbf{2. No semantic similarity:} Words are orthogonal one-hot vectors, so ``car'' and
``automobile'' have zero similarity. This causes vocabulary mismatch in retrieval.

\textbf{3. High dimensionality and sparsity:} With vocabularies of $10^5$--$10^6$ words,
each document vector is extremely sparse. Short documents are especially affected.

\textbf{4. No contextual disambiguation:} ``bank'' (financial vs. river) receives the
same feature weight regardless of context.

Subsequent methods addressed these systematically: \emph{N-grams} capture local order
(bigrams, trigrams). \emph{TF-IDF} addresses the frequency bias. \emph{Word2Vec and GloVe}
provide dense representations where similar words have similar vectors. \emph{RNNs and
LSTMs} model sequential dependencies. \emph{Attention mechanisms} enable direct pairwise
interaction between any two positions. \emph{Transformers and BERT} provide fully
contextualized representations where each token's embedding depends on the entire input.

\redflags
\begin{itemize}
    \item Only mentioning one problem (e.g., ``it loses word order'') without depth.
    \item Not connecting BoW limitations to the motivations for subsequent methods.
    \item Being unable to give a concrete failure example.
\end{itemize}

\interviewtest{Understanding of representation limitations and the progression of ideas
that led from BoW to transformers.}

\goodgreat
\begin{itemize}
    \item \badgeLfive{L5}: Names three or more problems with concrete examples, mentions at
    least one method that addresses each.
    \item \badgeLsix{L6}: Traces the full arc from BoW to transformers, explaining what
    each method added and what it still lacked.
    \item \badgeLseven{L7}: Discusses why BoW is \emph{still} useful (fast, interpretable,
    effective for certain tasks) despite its limitations, and when the added complexity of
    neural methods is not warranted.
\end{itemize}

\followups{When would BoW still be your first choice? How does TF-IDF address the
frequency bias but not the semantic similarity problem? What did Word2Vec add that
n-grams could not?}
\end{interviewq}

\vspace{0.5em}

% --- Question 7: Multilingual Preprocessing System Design ---
\begin{interviewq}
\textbf{Q: Design a text preprocessing pipeline for a multilingual sentiment analysis system covering 20+ languages.}

\badgeLsix{L6} \quad \badgetype{System Design} \quad \badgefreq{\freqmed}

\stronganswer
I would design this pipeline in layers, from language-agnostic to language-specific:

\textbf{Layer 1 -- Encoding and normalization:} Convert all input to UTF-8. Apply NFKC
Unicode normalization to handle equivalent character representations. This is universal
across languages.

\textbf{Layer 2 -- Language identification:} Use a fast LangID model (e.g., fastText's
language detector, which supports 170+ languages in microseconds) to route text to
language-specific processing.

\textbf{Layer 3 -- Language-aware tokenization:} This is the most critical decision.
Options include: (a) A multilingual subword tokenizer (SentencePiece trained on
multilingual data, as used in mBERT, XLM-R), which handles all languages uniformly. (b)
Language-specific tokenizers where needed (e.g., Jieba or pkuseg for Chinese word
segmentation, MeCab for Japanese). I would use option (a) for the model input and option
(b) for feature extraction in classical pipelines.

\textbf{Layer 4 -- Selective normalization:} Lowercasing is applied for most Latin-script
languages but skipped for case-sensitive languages and NER tasks. Stopword removal is
applied only if using classical models, with language-specific stopword lists. Stemming
and lemmatization are applied per-language where resources exist (Snowball for European
languages); for unsupported languages, the subword tokenizer handles morphological
variation.

\textbf{Layer 5 -- Special handling:} Emoji normalization (map to text descriptions or
preserve as features---emojis carry strong sentiment signal). URL and mention
normalization. Code-switching detection for text that mixes languages within a sentence.

\textbf{Design considerations:} The tokenizer tax---multilingual subword tokenizers produce
more tokens for non-Latin scripts, effectively reducing context window capacity.
Transliteration as an option for very low-resource languages. Separate sentiment lexicons
per language for feature-based approaches.

\redflags
\begin{itemize}
    \item Proposing an English-only pipeline and saying ``translate first.''
    \item Not mentioning language identification as a first step.
    \item Ignoring CJK segmentation challenges.
    \item Not considering the tokenizer tax for non-Latin scripts.
\end{itemize}

\interviewtest{Production NLP experience, awareness of multilingual challenges, ability to
design systems that work at global scale.}

\goodgreat
\begin{itemize}
    \item \badgeLfive{L5}: Mentions language detection, Unicode normalization, and
    language-specific tokenization.
    \item \badgeLsix{L6}: Designs a complete pipeline with clear layers, discusses
    tradeoffs (translate-then-analyze vs. multilingual model), and mentions the tokenizer
    tax.
    \item \badgeLseven{L7}: Discusses how to handle code-switching, transliteration for
    low-resource languages, evaluation challenges across languages (annotation quality
    varies by language), and data augmentation strategies (back-translation, cross-lingual
    transfer).
\end{itemize}

\followups{How would you handle a language not in your tokenizer's training data? What is
cross-lingual transfer and when does it work? How do you evaluate sentiment analysis
quality across languages with different annotation standards?}
\end{interviewq}

\vspace{0.5em}

% --- Question 8: Context in NLP ---
\begin{interviewq}
\textbf{Q: Why is context modeling crucial in NLP? Give examples where lack of context leads to failures.}

\badgeLfive{L5} \quad \badgetype{Conceptual} \quad \badgefreq{\freqmed}

\stronganswer
Context is crucial because natural language is fundamentally ambiguous, and context is what
resolves that ambiguity. Without context, NLP systems fail in systematic ways.

\textbf{Lexical ambiguity:} The word ``Apple'' could be a fruit or a company. Without
sentence context (``Apple released a new iPhone'' vs. ``I ate an apple''), any NLP system
will misclassify entities, topics, or sentiment. Static word embeddings like Word2Vec
assign a single vector to ``Apple'' regardless of context---this is precisely why
contextual embeddings (ELMo, BERT) were a breakthrough.

\textbf{Negation and sentiment:} ``This movie is not good'' has the opposite sentiment from
``This movie is good.'' BoW models that ignore word order conflate these. Even simple
sentiment lexicon approaches fail because ``not'' + ``good'' requires compositional
understanding, not just keyword matching.

\textbf{Coreference:} ``The doctor told the nurse she was wrong.'' Without discourse
context, ``she'' is ambiguous. Resolving it requires understanding the conversation history
or world knowledge.

\textbf{Pragmatics and implicature:} ``Can you open the window?'' is a request, not a
question about ability. Dialogue systems without pragmatic context modeling respond with
``Yes, I can'' instead of actually opening the window (or acknowledging the request).

\textbf{Cross-sentence dependencies:} In summarization and QA, the answer often depends
on information spread across multiple sentences or paragraphs. Models that process
sentences independently miss these connections.

The evolution from context-free representations (BoW, TF-IDF, Word2Vec) to contextual ones
(ELMo, BERT, GPT) is fundamentally about building models that can use context at
progressively larger scales---from local word windows to full document context to
multi-document context in RAG systems.

\redflags
\begin{itemize}
    \item Giving only one example and not connecting it to model design.
    \item Not mentioning the static-to-contextual embedding evolution.
    \item Confusing ``context'' with ``more data.''
\end{itemize}

\interviewtest{Fundamental understanding of why NLP is hard, ability to connect
linguistic challenges to architectural decisions.}

\goodgreat
\begin{itemize}
    \item \badgeLfive{L5}: Gives three or more concrete examples spanning different types
    of context (lexical, sentential, discourse) and connects to model choices.
    \item \badgeLsix{L6}: Traces the evolution from context-free to contextual
    representations, explains how transformers model context via self-attention, and
    discusses context window limitations.
    \item \badgeLseven{L7}: Discusses the frontier challenges---long-range context (``lost
    in the middle''), cross-document reasoning, knowledge that is not in any document
    (commonsense), and how RAG extends context beyond the training data.
\end{itemize}

\followups{How does BERT's attention mechanism model context? What happens when the
relevant context exceeds the model's context window? How does RAG address context
limitations?}
\end{interviewq}

\vspace{0.5em}

% --- Question 9: Perplexity ---
\begin{interviewq}
\textbf{Q: What is perplexity? How is it computed, and what are its limitations as an evaluation metric?}

\badgeLfive{L5} \quad \badgetype{Conceptual} \quad \badgefreq{\freqhigh}

\stronganswer
Perplexity is the standard intrinsic evaluation metric for language models. It measures how
well a probability model predicts a held-out test set. Formally, for a test sequence of
$T$ tokens:

\[
\text{PP}(W) = \exp\!\left(-\frac{1}{T}\sum_{t=1}^{T}\log P(w_t \given w_{<t})\right)
\]

This is the exponential of the cross-entropy. Intuitively, perplexity represents the
effective branching factor: a perplexity of 50 means the model is, on average, as uncertain
as if choosing uniformly among 50 options at each step. Lower is better.

\textbf{Limitations:} (1) \emph{Vocabulary dependence}---models with different tokenizers
produce incomparable perplexities because they predict different units. A character-level
model and a subword-level model cannot be directly compared. (2) \emph{Does not capture
quality}---perplexity measures prediction accuracy, not fluency, coherence, factuality, or
usefulness of generated text. A model could have low perplexity but generate toxic or
factually incorrect content. (3) \emph{Domain sensitivity}---perplexity measured on
Wikipedia does not predict performance on clinical notes. (4) \emph{Equal weighting}---all
tokens contribute equally, but predicting function words (``the,'' ``of'') is trivially
easy while predicting content words or rare terms is what matters for downstream tasks.

For LLM evaluation, downstream benchmarks (MMLU, HumanEval, GSM8K) and human evaluation are
far more informative than perplexity alone.

\redflags
\begin{itemize}
    \item Not knowing the formula or the connection to cross-entropy.
    \item Saying ``lower perplexity means the model is better at generating text'' without
    qualification.
    \item Not mentioning the vocabulary dependence problem.
\end{itemize}

\interviewtest{Understanding of language model evaluation, mathematical rigor, awareness
of metric limitations.}

\goodgreat
\begin{itemize}
    \item \badgeLfive{L5}: States the formula, gives the branching factor intuition, names
    two limitations.
    \item \badgeLsix{L6}: Explains the cross-entropy connection, discusses bits-per-character
    as an alternative for cross-tokenizer comparison, and connects to downstream evaluation.
    \item \badgeLseven{L7}: Discusses calibration (does low perplexity imply the model's
    probabilities are well-calibrated?), the relationship between perplexity and downstream
    performance (weak correlation for capable models), and how perplexity relates to
    scaling laws.
\end{itemize}

\followups{How would you compare perplexity between models with different tokenizers?
What is bits-per-byte? What is the relationship between perplexity and cross-entropy loss?}
\end{interviewq}

\vspace{0.5em}

% --- Question 10: Naive Bayes vs. BERT ---
\begin{interviewq}
\textbf{Q: When would you choose Naive Bayes over a fine-tuned BERT for text classification? Justify your answer.}

\badgeLfive{L5} \quad \badgetype{Conceptual} \quad \badgefreq{\freqmed}

\stronganswer
I would choose Naive Bayes over fine-tuned BERT in several specific scenarios:

\textbf{1. Very small training sets ($<$ 100 examples):} Naive Bayes has far fewer
parameters to estimate. With 50 labeled examples, BERT fine-tuning is prone to overfitting
on the classification head, while Naive Bayes with TF-IDF features can still produce
reasonable results. BERT's advantage comes from pretraining, but fine-tuning requires enough
task-specific data to adapt the classification head.

\textbf{2. Extreme latency requirements ($<$ 1 ms):} Naive Bayes inference is a sparse
vector multiplication---microseconds. BERT inference requires a forward pass through 12+
transformer layers---milliseconds to tens of milliseconds even with optimization. For
systems processing millions of requests per second (spam filtering, content routing), this
difference is significant.

\textbf{3. Interpretability requirements:} Naive Bayes feature weights directly show which
words drive the prediction. ``This email was classified as spam because the words `free,'
`winner,' and `click' have high spam likelihood ratios.'' BERT's predictions are opaque.

\textbf{4. Resource constraints:} Naive Bayes requires no GPU, minimal memory, and trains
in seconds. BERT fine-tuning requires a GPU and minutes to hours of training.

\textbf{5. Baseline establishment:} Naive Bayes is the first model I would train on any
text classification task. If it achieves 95\% accuracy, the marginal gain from BERT may not
justify the complexity.

That said, for most production text classification tasks with reasonable data ($>$ 1,000
examples) and moderate latency budgets ($>$ 10 ms), fine-tuned BERT or a distilled model
will significantly outperform Naive Bayes.

\redflags
\begin{itemize}
    \item Saying ``always use BERT because it is better'' without considering context.
    \item Saying ``Naive Bayes is always worse'' without acknowledging its advantages.
    \item Not quantifying the latency or data requirements.
\end{itemize}

\interviewtest{Practical judgment about model selection, understanding of the
simplicity-performance tradeoff, production awareness.}

\goodgreat
\begin{itemize}
    \item \badgeLfive{L5}: Gives at least two scenarios favoring Naive Bayes with clear
    reasoning.
    \item \badgeLsix{L6}: Discusses the Pareto frontier of accuracy vs. latency/cost,
    mentions distilled models (DistilBERT) as a middle ground, and discusses when the
    additional complexity of BERT is justified.
    \item \badgeLseven{L7}: Discusses the full model selection hierarchy for production
    (regex $\to$ rules $\to$ NB/SVM $\to$ distilled $\to$ BERT $\to$ LLM), cost-benefit
    analysis, and how to decide when to ``upgrade'' from a simple model.
\end{itemize}

\followups{What about DistilBERT or TinyBERT as middle-ground options? How would you
quantify the business value of the accuracy improvement from BERT? What about few-shot
classification with an LLM as an alternative?}
\end{interviewq}

\bigskip

\begin{keyinsight}[Chapter 1 Summary: What Interviewers Are Really Testing]
When interviewers ask about classical NLP, they are testing three things: (1) \emph{Do you
have depth?}---Can you explain the math behind TF-IDF, derive BM25, and reason about
smoothing in language models? (2) \emph{Do you have practical judgment?}---Do you know when
a simple model is the right choice? (3) \emph{Do you understand the arc?}---Can you explain
why each paradigm shift happened and what limitations drove it? Candidates who demonstrate
all three signal that they are not just users of NLP tools but engineers who understand the
foundations these tools are built on.
\end{keyinsight}
